\section{问题三：连续通信保障下的运输—中继联合调度}
\label{sec:q3}

\subsection{问题分析与求解思路}
\subsubsection{运输轨迹中的通信需求}
运输任务满足载荷、能量和箱级时限，并不自动保证飞行全过程能与地面网关通信。通信能力随三维位置和地形遮挡变化；中继同时需要满足运输机接入与网关回传，而不是只靠近运输机即可。无人机通信的部署与移动性具有相互影响\cite{zeng2016uavcom,zeng2019sky}，本题将其具体化为每条运输轨迹上的连续链路约束。

\subsubsection{中继位置、服务窗口与运输时刻的耦合}
一处悬停位置是否有效，取决于它能在何时保障哪些航段，以及中继实体和能源组件能否在该时刻到位。提前运输任务可能扩大中继接续压力，移动中继点也会改变转场时间、能耗和通信切换时刻。本文因此先计算候选位置对应的连续可服务区间，再联合选择服务任务与资源顺序，最后精修连续开始时刻。中继与运输均使用真实任务生命周期，不把静态覆盖当作已完成调度。

\subsection{连续通信与联合调度模型}
\subsubsection{直连和单中继链路可用条件}
设端点$a,b$在时刻$t$的三维距离为$d_{ab}(t)$。采用给定自由空间损耗关系\cite{iturp525}和地形附加损耗：
\begin{equation}
 L_{ab}(t)=32.45+20\log_{10}f+20\log_{10}\frac{d_{ab}(t)}{1000}
 +10\,\zeta_{ab}(t),
\label{eq:q3_loss}
\end{equation}
$f=2400$\,MHz，$\zeta_{ab}(t)$表示传播线是否被模型中的地形遮挡。设备功率、天线增益、系统损耗与接收要求共同给出双向允许损耗门限：运输机—网关122\,dB、运输机—中继116\,dB、中继—网关126\,dB。附加传播损耗为$\Delta$时，链路可用条件是
\begin{equation}
 \chi_{ab}(t)=\mathbf1\{L_{ab}(t)+\Delta\le L_{ab}^{\max}\}.
\label{eq:q3_link}
\end{equation}
遮挡计入损耗，而不是一律否决；是否可用仍取决于总损耗与门限。实际网关高度来自输入，不以O01地面位置代替其三维端点。

运输机$U$优先直连网关$G$。直连不可用时，一项在线中继$q$必须同时满足$U\leftrightarrow R_q$和$R_q\leftrightarrow G$。因此
\begin{equation}
 \chi_U(t)=\max\left\{\chi_{UG}(t),\max_{q:\,t\in[a_q,b_q]}
 \chi_{UR_q}(t)\chi_{R_qG}(t)\right\},\qquad \chi_U(t)=1.
\label{eq:q3_provider}
\end{equation}
每时刻只选择一条直连或单中继路径，不构造多级中继网络。提供方的选择需与其真实服务窗口对应；几何可达但尚未上线的中继不能计入式\eqref{eq:q3_provider}。

\subsubsection{链路距离门限与地形状态}
总损耗约束可以在地形状态固定时改写成距离门限。给定遮挡状态$\zeta\in\{0,1\}$和额外损耗$\Delta$，由式\eqref{eq:q3_loss}得
\begin{equation}
 d_{ab}^{max}(\zeta,\Delta)=1000\times10^{\{L_{ab}^{max}-32.45-20\log_{10}f-10\zeta-\Delta\}/20}.
 \label{eq:q3_range_threshold}
\end{equation}
距离单位为m，频率仍用MHz。同一设备下，10 dB遮挡使允许距离变为无遮挡的$10^{-1/2}$，约31.62\%。因而轨迹即使只移动很小距离，只要跨过遮挡事件，也可能改变通信可行性。一个与地形无关的圆形覆盖圈不足以描述这种关系。

固定遮挡状态时，附加0.25 dB使允许距离乘以$10^{-0.25/20}\approx0.9716$，约缩小2.84\%。这个变化可能移动关键服务区间端点，进而影响运输开始时刻。接入与回传又必须同时可用，所以把中继移近运输机改善一跳，并不保证另一跳同步改善。这些关系将空间位置、悬停高度和服务时间联系起来。

\subsubsection{阶段轨迹与通信事件分割}
在一段爬升、巡航或下降中，运输位置可写为$\boldsymbol x(\tau)=\boldsymbol x_0+\boldsymbol v\tau$。对固定通信端点$\boldsymbol p$，距离平方是
\begin{equation}
 d^2(\tau)=\|\boldsymbol v\|^2\tau^2+2\boldsymbol v^{\mathsf T}(\boldsymbol x_0-\boldsymbol p)\tau+\|\boldsymbol x_0-\boldsymbol p\|^2.
 \label{eq:q3_distance_quadratic}
\end{equation}
在遮挡状态不变的区间内，将此式与相应门限平方联立，可以确定进入或离开距离可行域的时刻。交接阶段的位置固定，则退化为常数判定。

地形遮挡部分沿传播线扫过的射线扇面求与闭合像元棱柱的接触事件，再与距离根、阶段端点和服务窗口端点合并排序。相邻事件间没有新状态变化，可用区间内部状态代表该开区间；每个端点另行检查，保留瞬时擦边事件。连续条件因此被转化为有限的区间与端点对象，而不是以稀疏采样近似宣称全过程连通。

\begin{table}[htbp]\centering\small
\caption{连续通信判定中需要保留的事件}\label{tab:comm_events}
\begin{tabularx}{\linewidth}{p{2.8cm}LL}\toprule
事件类型 & 来源 & 对计算的作用\\\midrule
阶段切换 & 爬升、巡航、下降与交接转换 & 更新位置函数和速度\\
地形接触 & 传播线开始或停止触及地形棱柱 & 更新遮挡损耗状态\\
距离门限 & 端点间距达到允许值 & 更新距离可行状态\\
服务窗口端点 & 中继建链完成或停止服务 & 更新在线提供方集合\\
瞬时擦边 & 边、角点和孤立接触 & 对单独时刻核验闭边界条件\\\bottomrule
\end{tabularx}\end{table}

Saunders等对自主配送无人机技术的综述同时讨论航行、通信与任务执行需求\cite{saunders2024}。本题将这些关系集中到可计算的轨迹与服务接口：几何模块说明某链路何时可用，联合排程说明提供方何时真正上线，最终验证检查它们能否覆盖运输执行的全部必要时段。


\subsubsection{中继转场、悬停与能源消耗}
设悬停点$p$的地面高程为$z_g(p)$，悬停海拔为$z_p$，离地高度满足$0<z_p-z_g(p)\le300$\,m。中继转场必须兼容悬停端点，取
\begin{equation}
 H_R(p)=\max\left\{\max_{c\in\mathcal C(O01,p)}Z_c+50,\ z_p,\ h_{O01}\right\}.
\label{eq:q3_relay_height}
\end{equation}
此规则只用于中继，普通运输仍遵守式\eqref{eq:common_cruise_altitude}。对于给定位置分别计算出程与回程时间$t_p^{\mathrm{out}},t_p^{\mathrm{back}}$和转场能耗$E_p^{\mathrm{tr}}$，不能抬高悬停点却保留原转场代价。

中继准备时间为180\,s，建链30\,s，返航周转300\,s；电池可用能量3.2\,kWh，悬停与通信功率合计1.10\,kW。任务$q$的开始为$s_q^R$、实际服务窗口为$[a_q,b_q]$，则
\begin{equation}
 a_q=s_q^R+180+t_{p_q}^{\mathrm{out}}+30,
 \qquad e_q^R=b_q+t_{p_q}^{\mathrm{back}},
\label{eq:q3_lifecycle}
\end{equation}
\begin{equation}
 E_q^R=E_{p_q}^{\mathrm{tr}}+\frac{1.10}{3600}(30+b_q-a_q),
 \qquad E_q^R\le0.8\times3.2.
\label{eq:q3_relay_energy}
\end{equation}
准备时间未额外添加题面未给定的耗电项。中继实体占用延续至$e_q^R+300$，能源组件延续至$e_q^R+\tau(SOC_q)$，两者分别检查，不能将周转、充电和最后返航混为一个时刻。

\subsubsection{连续服务、资源接续与联合完成时间}
对于给定运输任务$i$，把相对开始的直连失效时间集合划成需求区间$h=[\ell_h,u_h]$。记$\mathcal A_h$为能覆盖该整个区间且回传可用的候选中继任务，$z_{hq}$表示选择$q$保障$h$，$y_q$表示启用中继任务。典型覆盖约束为
\begin{equation}
 \sum_{q\in\mathcal A_h}z_{hq}=1,\quad z_{hq}\le y_q,\quad
 z_{hq}=1\Rightarrow
 \begin{cases}a_q\le s_i+\ell_h,\\ b_q\ge s_i+u_h.\end{cases}
\label{eq:q3_cover}
\end{equation}
候选区间的几何条件已确保其内部双跳链路可用；边界接触仍需单独核验。中继实体与能源组件均须唯一指派，同一实体的占用不能重叠；运输侧继续满足箱级时限、20\%返航余量及飞机—电池约束。在零迟到条件下最小化
\begin{equation}
 T_{\max}^{\mathrm{joint}}=\max\left\{\max_i(s_i+p_i),\max_q e_q^R\right\}.
\label{eq:q3_objective}
\end{equation}
最后返航后的充电和中继周转参与资源接续，但不额外延长式\eqref{eq:q3_objective}。时间优先、节能及增强通信保障分别保留方案，不将dB、kWh和秒任意加权为一个分数。

\subsubsection{服务窗口与两类中继资源的连接}
几何覆盖和时间覆盖分别承担不同约束。候选位置属于$\mathcal A_h$，表示能够满足需求区间$h$的接入和回传；实际指派$z_{hq}=1$后，还要求其在线窗口包含运输的绝对需求区间。几何条件成立而中继未到位时，不能向运输任务提供服务。

中继任务同时占用实体机和能源组件。若$q'$使用执行$q$的同一架中继机，需要等到$q$返航并完成周转；若使用同一能源组件，则需要等到该组件充满。两种接续分别为
\begin{equation}
 s_{q'}^R\ge e_q^R+t^{turn},\qquad
 s_{q'}^R\ge e_q^R+c_q^R,
 \label{eq:q3_two_relay_resources}
\end{equation}
各式仅在对应资源相同且$q$为前项时启用。题目允许同型资源复用，因此同一中继机可更换组件，已恢复组件也可供另一实体机使用。

服务结束$b_q$同时决定返航时刻、悬停能耗和能源恢复。延长窗口可能覆盖更多运输需求，却推迟后续资源；删去不必要的早到悬停，则可能同时节能并改善接续。将服务起止保留为连续变量，才能把通信保障真正纳入调度，而不是在运输时间表外另放若干静态点。

\subsubsection{分层计算中的决策接口}
候选重定位负责新的位置、高度和传播关系，联合整数模型负责任务选择、提供方和资源路径，固定结构线性规划负责最后的连续时刻。空间计算不直接塞进整数规划的代数约束，而是生成可审查的可服务关系；给定离散选择后，时间条件尽量交给线性优化精确处理。

这种接口还说明了扰动后的处理层次：只改变中继最早上线时刻，可复用几何并修复时序；传播门限改变，需要重检距离根与可服务区间；候选不再能覆盖关键需求，则回到提供方或位置选择层。每次调整都以同一运输需求和物理规则为基础，避免不同模块分别“优化”后得到互不兼容的日程。


\subsection{中继部署与任务时序的联合求解}
\subsubsection{候选悬停位置生成与可服务区间计算}
本节区分题目要求的联合决策空间与实际搜索的有限子域。求解从已保存的23项运输任务和4个中继位置起步，对位置、高度、资源顺序和通信切换进行局部改进；另以小规模运输任务重组检验是否能进一步改善。候选的初始箱组、路线及部分悬停位置继承自补充资料，具体来源与本项目新增计算分列于附录\ref{app:source_provenance}，不把所有初始结构描述为从零构造。

\begin{table}[htbp]\centering\small
\caption{联合求解各阶段的开放变量与固定范围}\label{tab:q3_decision_scope}
\begin{tabularx}{\linewidth}{p{2.4cm}XX}\toprule
阶段 & 允许改变的对象 & 固定内容与输出\\\midrule
位置候选搜索 & 单个中继的水平位置和离地高度 & 暂固定运输模式与其他中继；输出少量候选，尚不宣称完整可行\\
运输重组对照 & 小货箱并集中的组批、机型、访问顺序及选中任务数 & 只在显式给定的有限任务池内选择；未截断的枚举与截断池分开记录\\
联合整数排程 & 候选任务选择、开始时刻、提供方、飞机/电池及中继实体顺序 & 单个任务的路线、单个中继槽的位置已给定；输出可行离散结构\\
连续时间精修 & 运输开始、中继服务开始/结束及联合工期 & 固定路线、位置、提供方、真实资源顺序与充电分段；输出时刻\\
独立验收 & 不再改变计划 & 从完整几何重建区间、端点、能量与资源占用，接受或拒绝候选\\\bottomrule
\end{tabularx}
\end{table}

位置生成在当前点附近按水平网格偏移，并枚举25、50、\ldots、300\,m共12个离地高度。默认搜索半径1200\,m、网格300\,m、每个角色保留10个候选；保存的局部试验还使用75\,m和50\,m尺度，以及面向关键角色的扩大探测，均为有限试验域。先检查DEM覆盖、AGL、回传链路和中继转场，再以任务阶段严格内部、约25\,s间隔的采样估计接入潜力。评分按角色优先考虑出程、回程或二者之和，并以链路裕度破同分；保留水平位置互异的短名单。这里的抽样只可能筛选候选，不能作为整个连续空间不可行的证据。

每个保留点都重新建立完整几何。运输机在一个飞行阶段内位置随时间线性变化；固定端点至该阶段的射线扇面与DEM地形棱柱求交，得到遮挡事件区间，再加入距离门限的解析交点、飞行阶段和服务窗口边界。相邻端点之间的开区间内状态固定，检查内部状态并单独检查全部临界端点，保留孤立擦边事件。这一步才产生式\eqref{eq:q3_cover}使用的可服务关系，不用采样连线替代。

候选搜索先比较完整求解后的工期，再比较同工期能耗；当前实现的接受阈值为工期改善超过0.005\,s，或工期差不超过0.005\,s时能耗改善超过0.005\,kWh。该阈值仅用于有限搜索的数值筛选，不代表实际执行的抗扰动余量。一次角色扫描或预定网格完成后停止，不能宣称无导数全空间搜索已经收敛。

\subsubsection{通信提供方与资源顺序的联合选择}
对有限运输任务池$\mathcal I$取选择变量$x_i$，逐箱覆盖沿用式\eqref{eq:q2_cover}；只重排已有23项任务时固定$x_i=1$。每个中继槽$q$对应一个候选位置，取启用变量$y_q$及式\eqref{eq:q3_cover}的提供方变量$z_{hq}$。需求属于任务$i$时，将覆盖等式右端写为$x_i$，未选择任务不再产生通信需求；$z_{hq}\le y_q$，同时禁止启用不服务任何需求的中继槽。

有限模型通过\emph{任务可接续路径}控制资源数量。对同型飞机，二元变量$v^U_{ij}$表示任务$j$紧接$i$，$o_i^U$表示任务$i$占用一条资源链的起点，满足
\begin{equation}
\begin{aligned}
 o_i^U+\sum_jv^U_{ji}&=x_i,&\sum_jv^U_{ij}&\le x_i,\\
 \sum_{i:m_i=m}o_i^U&\le n_m^U,&
 v^U_{ij}=1&\Rightarrow s_j\ge s_i+p_i .
\end{aligned}
\label{eq:q3_transport_paths}
\end{equation}
电池路径将$p_i$替换为$p_i+c_i$，起点数上限替换为该型号电池库存。正时长的先后约束排除环，故每条路径可以赋予一个实体资源编号。中继实体同样构造路径，接续条件为后项任务开始不早于前项返航加300\,s周转。条件式以二元变量和由有限时间上界计算的大常数实现，不用一个无根据的无限大数。

选定位置后，中继服务开始$a_q$、结束$b_q$以及任务开始$\sigma_q$满足式\eqref{eq:q3_lifecycle}的启用形式：$\sigma_q=a_q-\ell_qy_q$，$\ell_q=180+t^{out}_{p_q}+30$；$b_q\ge a_q$，能耗为转场项、建链项及服务时长项之和。服务覆盖用$z_{hq}$连接运输和中继时间轴。运输硬截止在该阶段另保留30\,s的算法性余量，期望送达仍要求零迟到；30\,s不是新添题设，而是本次搜索采用的保守限制。

当前中继架次试验最多启用6个槽，初始满电能源组件为6组，因此内层整数模型可先让每个启用中继任务使用不同组件，保证存在不复用的组件指派，不需要在此阶段对组件再建组合接续变量。选出任务后，按真实含充电区间重新构造组件复用链，正式方案只需3组；后续连续精修显式保留这3组组件的恢复关系。这一简化依赖“启用任务数不超过初始组件数”，不能直接推广到更多中继架次。

先用HiGHS求有限模型的最小联合工期，再在已得工期增加不超过$10^{-5}$\,s的条件下，允许整数关系继续改变并最小化总能耗。接着固定实际资源链进入下一节的连续精修。当前保存的设置为单线程、整数相对间隙$10^{-7}$、整数可行性容差$10^{-7}$；公开重复优化入口默认90\,s预算，历史各局部试验的实际限时分别保存在状态文件中，不把入口默认值当作所有试验的实际耗时。

\textbf{运输结构和架次数的作用。}除固定23项任务外，程序也对小货箱并集重新生成模式：并集不超过11箱时枚举非空子集、可用机型及路线；模式超过250项时按单位箱数时间、单位箱数能耗和最迟启动余量交错截断。保存对照覆盖S006/S007、S004/S005、S003/S015等组合，联合任务池最大361项，未得到比对应参考更好的已验证运输结构。因此本章主方案与同模型改进前方案的23项箱组、机型和有序路线逐项相同，约87.998\,s的成功改进应归于中继位置与联合时序，而不是运输重组。

\begin{table}[htbp]\centering\small
\caption{不同运输及中继架次试验的实际范围}\label{tab:q3_sortie_scope}
\begin{tabularx}{\linewidth}{p{3.0cm}XX}\toprule
对照 & 搜索范围与结果 & 可作出的结论\\\midrule
运输模式重组 & 小并集与有限路线池，最高361个候选；部分试验限时，未选出更好的正式结构 & 运输架次数可在池内改变；主方案23架次是最终选中结果，不是全局最少证明\\
3个中继架次 & 给定所测4个位置和完成时间上限，有限模型不可行 & 只排除该布局及时间阈值下的三架次方案\\
最多5或6架次 & 8个同位置重复出动槽；普通模型限时未改善 & 未改善或超时不证明所有增加架次的方案无效\\
最多5架次加对称消除 & 同位置槽安全重标号后，有限模型最优仍选4项 & 该有限模型无需第5项；不构成连续位置空间的四架次下界\\\bottomrule
\end{tabularx}
\end{table}

阶段性保存工期依次出现5924.967、5895.017、5866.141、5857.090和5836.970\,s。位置、资源顺序和能耗再优化在不同阶段同时调整，所以该轨迹是实际搜索进展，不是逐因素可加的消融分解。全部保存状态包括110份整数模型结果（91份有限域最优、15份有限域不可行、4份限时）及125份连续精修记录；记录数不能等同于不同有效方案数。

\subsubsection{给定任务结构后的时间安排优化}
\label{sec:q3_fixed_lp}
为使时刻精修可直接复核，本节给出实际使用的固定结构线性模型。输入固定运输箱组、机型、路线，中继位置，通信区间的提供方，以及飞机、电池、中继实体和组件的前后顺序。变量为运输开始$s_i$、中继服务开始$a_q$、服务结束$b_q$和联合最后返航$M$。记$\ell_q=180+t_{p_q}^{out}+30$，则任务开始$\sigma_q=a_q-\ell_q$、返航$e_q=b_q+t_{p_q}^{back}$。

对运输任务，根据箱级要求计算
\begin{equation}
 \bar s_i=\min\left\{\min_{b\in\mathcal B_i}(d_b-\delta_{ib}),\;
 \min_{b\in\mathcal B_i\cap\mathcal B_H}(h_b-\delta_{ib}-30)\right\},
 \quad 0\le s_i\le\bar s_i .
 \label{eq:q3_lp_start_bounds}
\end{equation}
$\mathcal B_H$是存在硬截止的货箱集合，空集最小值取正无穷。已固定给中继$q$的相对需求区间$[\ell_h,u_h]$满足
\begin{equation}
 a_q\le s_i+\ell_h-g_s,\qquad b_q\ge s_i+u_h+g_s,
 \quad a_q\ge\ell_q,\quad b_q\ge a_q .
 \label{eq:q3_lp_coverage}
\end{equation}
实际$g_s=10^{-4}$\,s是数值保护间隔。提供方在进入此模型前按完整可服务关系确定，线性规划本身不再自由改派或移动中继。

固定位置的转场能耗为$E_q^{tr}$，令$w=1.10/3600$\,kWh/s，则
\begin{equation}
 E_q=E_q^{tr}+w(30+b_q-a_q),\qquad 0\le E_q\le0.8B_R,
 \quad B_R=3.2\ \mathrm{kWh}.
 \label{eq:q3_lp_energy}
\end{equation}
充电按名义任务所在的分段固定，并同时施加保持该分段的能量条件，不能只线性化而遗漏分段边界。若$E_q\le0.1B_R$，则$c_q=(3.5\tau_R^{full}/B_R)E_q$；若$E_q\ge0.1B_R$，则
\begin{equation}
 c_q=\frac{0.65\tau_R^{full}}{0.9B_R}E_q
       +\tau_R^{full}\left(0.35-\frac{0.65\times0.1}{0.9}\right).
 \label{eq:q3_lp_charge}
\end{equation}
两段在$E_q=0.1B_R$处连续。该表达由式\eqref{eq:common_recharge}代入$SOC=1-E_q/B_R$得到。

对固定资源链上的相邻任务，分别要求
\begin{equation}
\begin{aligned}
 s_j&\ge s_i+p_i+g_r &&\text{（同一运输机）},\\
 s_j&\ge s_i+p_i+c_i+g_r &&\text{（同一运输电池）},\\
 a_v-\ell_v&\ge b_q+t_{p_q}^{back}+300+g_r &&\text{（同一中继机）},\\
 a_v-\ell_v&\ge b_q+t_{p_q}^{back}+c_q+g_r &&\text{（同一中继组件）}.
\end{aligned}
\label{eq:q3_lp_resource}
\end{equation}
其中$g_r=10^{-4}$\,s，运输$c_i$因任务属性固定而为常数；中继$c_q$由式\eqref{eq:q3_lp_energy}--\eqref{eq:q3_lp_charge}给出仿射函数。目标为
\begin{equation}
 \min M,\qquad s_i+p_i\le M,\quad b_q+t_{p_q}^{back}\le M .
 \label{eq:q3_lp_objective}
\end{equation}
还保留实现中的20000\,s变量上界；当前解远未触及该数值边界。求出$M_{LP}^*$后，再加$M\le M_{LP}^*+10^{-6}$\,s并最小化$\sum_q E_q$。两轮均使用SciPy HiGHS，原始及对偶可行性容差为$10^{-8}$。

该线性规划的工期下界为5836.969928328251\,s，二次节能后的保存值比其大$10^{-6}$\,s。因此可称给定路线、位置、提供方、资源顺序和充电分段的连续时序已闭合；不能称原问题全局最优，也不能把数值保护间隔或求解精度当作实飞时序容错。最终计划仍返回未修改的完整物理与连续通信验证器验收。


\subsubsection{给定结构的连续精修与复核接口}
运输路线和载荷固定后，运输作业与充电时长为常数；中继位置固定后，转场时间为常数，服务时长决定线性悬停能耗。再固定各组件处于哪一充电分段，并保留该段能量上下界，资源接续就具有仿射形式。这些条件共同解释固定结构模型为什么可以采用线性规划。

分段边界是模型的一部分。只使用一段充电直线而不约束能量范围，可能算出原函数不支持的恢复时刻。两轮优化先压缩联合工期，再在给定工期容许量内减少不必要的服务能耗，所得解仍须返回完整物理计算和通信事件核验。Huangfu与Hall给出的对偶修正单纯形研究为所用线性规划技术提供方法背景\cite{huangfu2018}；本文的固定结构结论则来自明确的输入、约束及保存下界。

位置和任务组合选择、资源先后安排、具体开始时刻分属不同规模的决策。分层方法使每一层都有可列出的输入与输出，也使恢复计算能只调整受扰动的层次。数值保护间隔用于提高计算稳定性，执行保障范围则由本章实际压力情景和连续核验单独评价。


\subsection{联合方案与结果验证}
\subsubsection{运输日程与中继部署方案}
工期优先方案由23个运输架次和4个中继架次组成，使用8架运输机、14组运输电池、2架中继机、3组中继能源组件。联合最后返航5836.969929\,s，总能耗68.939416\,kWh，其中运输66.251280\,kWh、中继2.688136\,kWh。80箱唯一交付，零加权迟到、零硬迟到，最小硬截止余量360.866938\,s。

全部23项运输任务的飞机、电池、开始、末箱交付与返航时刻见附录表\ref{tab:q3_full_schedule}，不沿用问题二的时刻；表\ref{tab:q3_sites}和表\ref{tab:q3_relays}给出四项中继任务的实际部署与服务时刻；编号“中继任务”与“实体中继机”严格区分。第一架实体先执行Q3-R-02，再执行Q3-R-01；第二架先执行Q3-R-03，再执行Q3-R-04。四个架次不是四架实体机。
\begin{table}[htbp]\centering\small
\caption{中继任务、实体资源与悬停位置}\label{tab:q3_sites}
\begin{tabular}{@{}lllrrr@{}}
\toprule
中继任务 & 实体 & 能源组件 & 经度 & 纬度 & AGL/m\\
\midrule
Q3-R-01 & R01 & R-ENG02 & 109.201003 & 23.033703 & 300\\
Q3-R-02 & R01 & R-ENG01 & 109.268309 & 23.014741 & 275\\
Q3-R-03 & R02 & R-ENG02 & 109.211343 & 23.020248 & 300\\
Q3-R-04 & R02 & R-ENG03 & 109.237677 & 23.059303 & 300\\
\bottomrule
\end{tabular}
\end{table}

\begin{table}[htbp]\centering\small
\caption{四项中继任务的实际服务时间表}\label{tab:q3_relays}
\begin{tabular}{@{}lrrrrr@{}}
\toprule
任务 & 开始/min & 服务开始 & 服务结束 & 返航/min & SOC/\%\\
\midrule
Q3-R-01 & 59.312 & 69.709 & 89.566 & 97.228 & 81.294\\
Q3-R-02 & 2.661 & 12.595 & 47.178 & 54.312 & 73.307\\
Q3-R-03 & 5.531 & 13.606 & 24.386 & 29.603 & 89.079\\
Q3-R-04 & 34.603 & 46.546 & 78.701 & 87.860 & 72.316\\
\bottomrule
\end{tabular}
\end{table}

\fig[.89]{q3_deployment.pdf}{工期优先方案的中继悬停位置与关键运输路线。中继点为实际位置，不用圆形覆盖圈代替地形条件下的可服务区域。}{fig:q3_deployment}
\fig[.97]{q3_resources.pdf}{四项中继任务的实体复用与服务窗口。深色表示实际服务；任务开始、服务开始、返航和周转分别处理。}{fig:q3_resources}

\subsubsection{连续通信及资源占用核验}
本轮从完整运行包重新构造名义几何和计划，八项核心导出与保存结果逐字节一致。独立验证完成17756项物理、能源、资源、完整区间和端点检查，失败0；另有36672个包含边界的时空核验点，失败0。保存的700份区间材料与1238份临界端点材料用于追溯覆盖，不能把核验计数当作实飞次数或随机稳定性样本量。

最低运输返航SOC为20.097383\%，最低中继SOC为72.315816\%。中继能源较宽裕不意味着通信时间宽裕。图\ref{fig:q3_provider}以关键任务的实际通信提供方显示切换关系，服务端点必须与中继在线窗口相互兼容。若只画静态连接网络，将无法发现这种时间瓶颈。
\fig[.93]{q3_providers.pdf}{关键运输任务的阶段与实际通信提供方。横轴为绝对调度时间，提供方切换来自正式区间记录，不是抽样插值。}{fig:q3_provider}

\subsubsection{运输与中继能源成本的分解}
联合方案运输能耗66.251280 kWh，中继能耗2.688136 kWh，中继占总能耗约3.90\%。表\ref{tab:rev_q3_transport_energy}按51条实际运输航段聚合，水平部分63.050347 kWh，爬升部分3.200933 kWh。爬升项在总量中所占比例较小，但仍会改变接近安全边界任务的机型和箱组可行性，因此在逐段计算中完整保留。
\begin{table}[htbp]\centering\small
\caption{联合方案运输能耗的物理分解}\label{tab:rev_q3_transport_energy}
\begin{tabular}{lrrr}\toprule
机型 & 水平能耗/kWh & 爬升能耗/kWh & 总能耗/kWh\\\midrule
A & 24.888150 & 1.516529 & 26.404679\\
B & 14.308489 & 0.814536 & 15.123025\\
C & 23.853709 & 0.869867 & 24.723576\\
\bottomrule\end{tabular}
\end{table}


四项中继任务的转场合计0.866261 kWh，建链和实际服务合计1.821876 kWh，后者约占中继能耗的67.77\%。表\ref{tab:rev_q3_relay_energy}显示，各任务服务10.780—34.583 min，占自身准备至返航作业时间的44.78\%—66.96\%。中继成本不仅取决于飞到哪里，还取决于在那里等待和服务多久。
\begin{table}[htbp]\centering\small
\caption{中继任务的转场、建链与悬停服务成本}\label{tab:rev_q3_relay_energy}
\begin{tabular}{lrrrrr}\toprule
中继任务 & 转场/kWh & 建链及服务/kWh & 总能耗/kWh & 服务/min & 服务占用比/\%\\\midrule
Q3-R-01 & 0.225372 & 0.373222 & 0.598594 & 19.858 & 52.37\\
Q3-R-02 & 0.210978 & 0.643196 & 0.854174 & 34.583 & 66.96\\
Q3-R-03 & 0.142681 & 0.206793 & 0.349474 & 10.780 & 44.78\\
Q3-R-04 & 0.287229 & 0.598665 & 0.885894 & 32.154 & 60.38\\
\bottomrule\end{tabular}
\end{table}


能源充足和服务及时是两种条件。中继返航SOC较高，提供了能源上的余量；但如果它未在关键切换前完成建链，运输任务仍可能等待。因而改善一个方案应先识别起作用的约束：位置可能主要改善遮挡和交接，开始时刻可能主要减少提前悬停，组件配置则主要影响跨架次恢复。这些作用都能在分解表与实际服务窗口中对应。

\subsubsection{通信提供方的实际保障工作量}
从正式保存的181个通信阶段、228条实际提供方区间重新聚合，运输任务累计需要34312.495780 s保障。固定网关承担20734.854799 s，占60.43\%；中继承担13577.640980 s，占39.57\%。18个运输架次在部分阶段使用中继，5个全程直连，见表\ref{tab:rev_q3_provider}。
\begin{table}[htbp]\centering\small
\caption{实际通信提供方的累计保障工作量}\label{tab:rev_q3_provider}
\begin{tabular}{lrrrr}\toprule
提供方 & 区间数 & 运输任务数 & 累计保障/min & 比例/\%\\\midrule
G01 & 112 & 23 & 345.581 & 60.43\\
Q3-R-01 & 41 & 7 & 92.738 & 16.22\\
Q3-R-02 & 28 & 7 & 48.093 & 8.41\\
Q3-R-03 & 15 & 3 & 24.757 & 4.33\\
Q3-R-04 & 32 & 9 & 60.706 & 10.62\\
\bottomrule\end{tabular}
\end{table}


累计保障时长按运输任务分别求和。如果两架运输机在同一秒使用同一中继，该秒贡献两份保障工作量。因此Q3-R-01约19.858 min的在线窗口，可对应约92.738 min的多任务累计保障。该指标说明其承担的任务服务量，不是系统失联时长、通信吞吐量或中继自己的飞行时长。当前模型没有额外的并发用户容量约束，保障关系按给定双跳链路条件核算。

实际提供方序列在合并相邻同提供方记录后共有47次切换，包括直连—中继和中继—中继交接。表\ref{tab:rev_q3_comm_tasks}列出中继保障需求较大的任务。相对于静态“覆盖多少服务区”，这些记录更直接回答：中继为哪些运输机服务、在什么时段服务、哪些切换会影响任务开始。
\begin{table}[htbp]\centering\small
\caption{中继保障需求较大的运输任务}\label{tab:rev_q3_comm_tasks}
\begin{tabular}{lrrr}\toprule
运输任务 & 直连/min & 中继/min & 提供方切换数\\\midrule
Q3-T-09 & 13.587 & 24.505 & 3\\
Q3-T-16 & 14.928 & 21.465 & 3\\
Q3-T-17 & 16.300 & 20.589 & 3\\
Q3-T-12 & 11.503 & 16.227 & 3\\
Q3-T-20 & 12.676 & 13.521 & 2\\
Q3-T-10 & 12.263 & 13.283 & 3\\
Q3-T-21 & 12.676 & 13.021 & 2\\
Q3-T-15 & 21.829 & 12.739 & 4\\
\bottomrule\end{tabular}
\end{table}


中继以较小的能源占比支撑较大比例的运输通信任务时间，体现了运输与通信资源的分工价值。该统计不替代优化证明，而为联合方案增加运行机理说明：空间部署提供可用链路，服务窗口保证时间匹配，实体和组件复用保证这些服务能够实际安排。


\subsubsection{通信切换对最后返航时间的影响}
与采用同一悬停兼容物理模型的改进前方案相比，当前方案工期缩短87.997560\,s。关键C型任务前段保障结束的相对时刻由956.973974延长至1039.150288\,s，即交接边界后移82.176314\,s；后继中继又提前5.821245\,s就绪，二者共同放松其最早开始约束。这些差值是同物理模型的保存对照，不能与更早的受限高度模型混合归因。

当切换约束活跃时，其形式可写为$s_i\ge a_q-\ell_h+g$，其中$g$是计算保护间隔。后继可服务时刻$a_q$提前、前段可持续时长$\ell_h$增加，都会使运输任务更早开始。当前末尾的S007—S003 C型任务与S011 A型任务同在5836.969929\,s返航；最后中继在5833.697980\,s返回。由此，优化中继的价值体现在减少运输等待，而不只是减少中继自己的飞行距离。

\subsubsection{工期优先与节能备选方案比较}
节能备选在5911.548秒完成，能耗为68.735791 kWh；比工期优先方案晚74.578秒、节省0.203625\,kWh。面向三组独立配置的折中方案为5866.140527\,s、68.975497\,kWh；它在时间与能耗上并不优于工期优先方案，但后续三组资源缺口较小，见第\ref{sec:q4}章。上述都是已验证的离散候选，不代表完整Pareto前沿或全局最优解。

\subsection{稳健性与敏感性分析}
\label{sec:q3_robust}
\subsubsection{附加传播损耗下的通信可行性}
分别固定工期优先方案与附加0.25\,dB保障备选的路线、位置、运输时刻和原中继提供方，在$\Delta=0,0.05,0.10,0.15,0.20,0.25$\,dB下重做完整区间与端点核验。允许遵守直连优先规则使用可用直连，但不切换到另一项未指派中继，也不移动日程。新增损耗门限解析交点与全部几何擦边端点一同检查，不能用原来的少量采样点判断新阈值。

表\ref{tab:q3_loss}表明：工期优先方案在名义条件通过，附加0.05\,dB即失去连续通信，按不同运输任务分别累计的失联时长为138.973270\,s；0.25\,dB时为503.429111\,s。该总量是各任务失联时长之和，不是整个系统墙钟时间的并集，也不是丢包量。增强保障备选在全部六个水平的区间和端点检查均通过。固定物理与提供方时，门限条件对$\Delta$单调，故通过0.25\,dB也支持该模型下整个$[0,0.25]$区间，不支持更大的未测损耗。
\begin{table}[htbp]\centering\small
\caption{附加损耗下两份固定方案的连续通信核验}\label{tab:q3_loss}
\begin{tabular}{@{}rrcrc@{}}
\toprule
损耗/dB & 原方案失联/s & 核验 & 增强方案失联/s & 核验\\
\midrule
0.00 & 0.000 & 通过 & 0.000 & 通过\\
0.05 & 138.973 & 失效 & 0.000 & 通过\\
0.10 & 300.470 & 失效 & 0.000 & 通过\\
0.15 & 368.384 & 失效 & 0.000 & 通过\\
0.20 & 436.035 & 失效 & 0.000 & 通过\\
0.25 & 503.429 & 失效 & 0.000 & 通过\\
\bottomrule
\end{tabular}
\end{table}

\fig[.88]{q3_loss.pdf}{固定两份方案后的附加损耗响应。纵轴为逐运输任务累计失联时长；增强保障备选与原工期优先方案分别核验，不混合统计。}{fig:q3_loss}

0.05\,dB是首个已测失效水平，不是精确最大容忍损耗。这里“失去通信”指当前模型和原提供方约定下的结果，不是实测链路失效概率。旧候选域的0.50/0.55\,dB结论不进入本试验。更重要的是，增强方案的成功不能被写成原工期优先方案本身能够承受0.25\,dB。

\subsubsection{中继到位延迟下的服务连续性}
对四项中继任务分别将最早可服务时刻向后移0、1、5、10、30、60\,s，运输时刻、原中继提供方及原定服务结束保持不变，形成24个情景。这模拟部署端迟发或上线延后造成的可用性丧失；转场路径与速度不变，不把它解释为飞行变慢或额外空中等待，因此不虚构额外飞行功率。

四项任务在所测1\,s延迟下均失效，但这不是精确临界延迟。以30\,s为例，Q3-R-01至Q3-R-04分别造成60.000、30.000、70.557、89.984\,s逐任务累计失联，说明被保障运输任务数和切换边界决定传播影响，不是延迟秒数的简单一比一映射。

随后对每项中继的5、30、60\,s延迟开放时序修复：路线、悬停位置、资源身份与顺序、通信覆盖指派和充电分段保持不变，只允许运输开始和中继服务边界向后调整，不提前任务。连续线性模型先最小化联合返航，再在该工期上尽量节省悬停能耗；12份输出全部回到完整物理模型重算，并经过区间、端点与独立时空点检查，零失败、零迟到。
\begin{table}[htbp]\centering\small
\caption{固定结构修复中继上线延迟的额外工期}\label{tab:q3_repair}
\begin{tabular}{@{}lrrrc@{}}
\toprule
延迟任务 & 延迟5秒 & 延迟30秒 & 延迟60秒 & 完整核验\\
\midrule
Q3-R-01 & 5.000 & 30.000 & 60.000 & 3/3通过\\
Q3-R-02 & 5.000 & 30.000 & 60.000 & 3/3通过\\
Q3-R-03 & 1.403 & 26.403 & 56.403 & 3/3通过\\
Q3-R-04 & 1.728 & 26.728 & 56.728 & 3/3通过\\
\bottomrule
\end{tabular}
\par\smallskip\parbox{.96\linewidth}{\footnotesize 中间三列均为新增工期/s，不是修复后总工期；转场物理参数保持不变。}
\end{table}

\fig[.89]{q3_repair.pdf}{给定任务结构下修复中继上线延迟的工期代价。每个点均经完整模型核验；曲线只连接已计算情景，不表示全局最优恢复策略。}{fig:q3_repair}

30\,s延迟的新增工期为26.403—30.000\,s，60\,s延迟为56.403—60.000\,s。部分修复后的中继能耗略减，因为延后部署缩短了无必要的提前悬停，并非飞行延迟可以“免费节能”。若实际情景包含额外空中等待，需要另设相应时间和功率再核验。上述结果支持存在时序恢复能力，但不能把恢复后的成功归为固定原方案稳健。

\subsubsection{保障裕度与工期、能耗的权衡}
附加0.25\,dB保障备选的联合工期5866.140527\,s，总能耗69.031751\,kWh，相比工期优先方案增加29.170597\,s和0.092335\,kWh。这是两份具体方案间的保障代价，不是所有增强保障方案的最小必要代价。针对上线延迟的12次修复也只在给定结构内优化，不扩张为自由位置、路径、分组同时可调的全局结果。

因此本问的稳健性结论是有条件的：名义最快候选在小幅传播损耗或1\,s上线延迟下已可能失效；增强损耗备选和明确的时序修复提供了可验证的补救。选择方案时应区分所需保障类型，不能用中继较高剩余电量或求解器微小间隙替代通信时间裕度。

\subsubsection{通信交接约束的恢复解释}
对关键运输任务$i$，若前段提供方可持续至相对时刻$\tau_i^{sw}$，后继中继$q$在绝对时刻$a_q$可用，含保护间隔$g$的交接要求可写为
\begin{equation}
 s_i+\tau_i^{sw}\ge a_q+g
 \quad\Longleftrightarrow\quad
 s_i\ge a_q-\tau_i^{sw}+g.
 \label{eq:q3_handoff_start}
\end{equation}
后继中继提前和前段保障延长，均可降低通信给任务开始施加的下界。当前同模型改进把这两方面结合，解释了约87.998 s工期下降。改进后两项运输末任务同刻返回，也说明下一轮位置调整可能遭遇新的资源瓶颈，收益不会无限线性叠加。

名义工期优先、增强通信保障和扰动后修复构成三类执行选择。附加0.25 dB保障备选比主方案增加约0.50\%工期、0.13\%能源；单项上线延迟则优先保持位置与路线、重新安排服务和运输时刻。12份修复给出了所测5、30、60 s上线延迟下的实际恢复方案，各自检查物理、时限、资源和连续通信。

\begin{table}[htbp]\centering\small
\caption{不同通信变化对应的计算与调整层次}\label{tab:q3_recovery_policy}
\begin{tabularx}{\linewidth}{p{2.3cm}LLL}\toprule
情景 & 处理方法 & 主要固定对象 & 交付结果\\\midrule
名义执行 & 使用工期优先方案 & 已验证任务、位置与资源结构 & 基准完整日程\\
持续附加损耗 & 对应保障备选与门限核验 & 所选备选的位置和任务 & 指定损耗条件下的日程\\
中继上线延后 & 连续时间修复 & 路线、位置、提供方与资源顺序 & 修复日程和额外代价\\
原覆盖关系失效 & 重新选择候选与任务 & 统一物理规则及实际需求 & 新结构经完整核验后使用\\\bottomrule
\end{tabularx}\end{table}

前三类已有保存方案或压力试验支持；覆盖结构重选是进一步扩展的求解入口。将扰动识别与调整权限明确对应，可以优先采用改动范围最小的恢复方式，同时在需要时返回上一决策层。通信保障由此体现为可选择的执行方案，而不止一个名义验证通过的标签。

